\section{导读：为什么一家 ToB 公司史也是 AI 产业课}

本节先说明这期为什么值得放进广义 AI/互联网队列。表面上，EP116 是吴明辉在上市前回顾明略科技、秒针系统、并购、现金流和 19 年 ToB 公司史；更深一层，它讨论的是企业级 AI 的真实落地条件：私有数据、行业 workflow、可审计交付、供给侧能力、Agentic Model、Tool Use 和人机互信。

这期最重要的不是“某家公司如何熬到 IPO”，而是一个反互联网爽文的 ToB 样本。ToC 互联网讲流量、增长和网络效应；ToB 企业服务讲客户现场、数据入口、交付质量、现金流、组织耐力和长期信用。AI 到来后，这些老问题不会消失，而是被重新排列：公开数据训练基础模型并卖 token 会越来越卷，拥有私有数据和行业流程的公司，才可能在企业级 Agentic Model 上产生差异化。

\begin{figure}[H]
\centering
\includegraphics[width=0.96\textwidth]{company-history-map.png}
\caption{明略 19 年公司史：第一段创业、第二段创业、痛苦急转和企业级 AI 重新连成一条线。自制概念图，依据 00:02:11--03:47:40 对谈内容整理。}
\end{figure}

\begin{knowledgebox}{读图：这不是线性创业史}
第一段是秒针、推荐系统、广告监测和 Web2 技术复杂度；第二段是明略、Palantir 式数据分析、ToB/ToG 路线和并购；痛苦急转是现金流、聚焦和经营能力；AI 时代则把这些积累重新解释为私有数据、企业 workflow 和 Agentic Model。
\end{knowledgebox}

\begin{importantbox}{本期核心命题}
企业级 AI 的关键不是“把通用模型卖给企业”，而是把模型接入企业私有数据、工具系统、业务流程和责任边界，让 AI 在真实工作中产生可审计、可交付、可计价的结果。
\end{importantbox}

\begin{knowledgebox}{本期阅读路线}
\begin{tabularx}{\textwidth}{p{0.22\textwidth}p{0.32\textwidth}X}
\toprule
路线 & 主要内容 & 要带走的问题 \\
\midrule
公司史 & 秒针、明略、并购、痛苦急转 & 为什么 ToB 公司需要长期积累，而不是一次增长爆发？ \\
数据线 & Web2 行为数据、conversational data、私有 data & 哪些数据能变成企业 AI 的差异化？ \\
产品线 & DeepMiner、Agentic Model、Tool Use & AI 怎样从分析工具变成工作流执行系统？ \\
组织线 & 供给侧能力、老板/合伙人、可信 AI & AI 会怎样重写企业组织和责任边界？ \\
\bottomrule
\end{tabularx}
\end{knowledgebox}

\subsection{本章小结}

EP116 的阅读方法是把公司史当作技术商业案例。前半段解释数据入口和 ToB 经营为什么难，后半段解释这些难点在 AI 时代为什么反而变成壁垒。

